% Options for packages loaded elsewhere
% Options for packages loaded elsewhere
\PassOptionsToPackage{unicode}{hyperref}
\PassOptionsToPackage{hyphens}{url}
\PassOptionsToPackage{dvipsnames,svgnames,x11names}{xcolor}
%
\documentclass[
  french,
  letterpaper,
  DIV=11,
  numbers=noendperiod,
  justified]{scrartcl}
\usepackage{xcolor}
\usepackage{amsmath,amssymb}
\setcounter{secnumdepth}{-\maxdimen} % remove section numbering
\usepackage{iftex}
\ifPDFTeX
  \usepackage[T1]{fontenc}
  \usepackage[utf8]{inputenc}
  \usepackage{textcomp} % provide euro and other symbols
\else % if luatex or xetex
  \usepackage{unicode-math} % this also loads fontspec
  \defaultfontfeatures{Scale=MatchLowercase}
  \defaultfontfeatures[\rmfamily]{Ligatures=TeX,Scale=1}
\fi
\usepackage{lmodern}
\ifPDFTeX\else
  % xetex/luatex font selection
\fi
% Use upquote if available, for straight quotes in verbatim environments
\IfFileExists{upquote.sty}{\usepackage{upquote}}{}
\IfFileExists{microtype.sty}{% use microtype if available
  \usepackage[]{microtype}
  \UseMicrotypeSet[protrusion]{basicmath} % disable protrusion for tt fonts
}{}
\makeatletter
\@ifundefined{KOMAClassName}{% if non-KOMA class
  \IfFileExists{parskip.sty}{%
    \usepackage{parskip}
  }{% else
    \setlength{\parindent}{0pt}
    \setlength{\parskip}{6pt plus 2pt minus 1pt}}
}{% if KOMA class
  \KOMAoptions{parskip=half}}
\makeatother
% Make \paragraph and \subparagraph free-standing
\makeatletter
\ifx\paragraph\undefined\else
  \let\oldparagraph\paragraph
  \renewcommand{\paragraph}{
    \@ifstar
      \xxxParagraphStar
      \xxxParagraphNoStar
  }
  \newcommand{\xxxParagraphStar}[1]{\oldparagraph*{#1}\mbox{}}
  \newcommand{\xxxParagraphNoStar}[1]{\oldparagraph{#1}\mbox{}}
\fi
\ifx\subparagraph\undefined\else
  \let\oldsubparagraph\subparagraph
  \renewcommand{\subparagraph}{
    \@ifstar
      \xxxSubParagraphStar
      \xxxSubParagraphNoStar
  }
  \newcommand{\xxxSubParagraphStar}[1]{\oldsubparagraph*{#1}\mbox{}}
  \newcommand{\xxxSubParagraphNoStar}[1]{\oldsubparagraph{#1}\mbox{}}
\fi
\makeatother


\usepackage{longtable,booktabs,array}
\newcounter{none} % for unnumbered tables
\usepackage{calc} % for calculating minipage widths
% Correct order of tables after \paragraph or \subparagraph
\usepackage{etoolbox}
\makeatletter
\patchcmd\longtable{\par}{\if@noskipsec\mbox{}\fi\par}{}{}
\makeatother
% Allow footnotes in longtable head/foot
\IfFileExists{footnotehyper.sty}{\usepackage{footnotehyper}}{\usepackage{footnote}}
\makesavenoteenv{longtable}
\usepackage{graphicx}
\makeatletter
\newsavebox\pandoc@box
\newcommand*\pandocbounded[1]{% scales image to fit in text height/width
  \sbox\pandoc@box{#1}%
  \Gscale@div\@tempa{\textheight}{\dimexpr\ht\pandoc@box+\dp\pandoc@box\relax}%
  \Gscale@div\@tempb{\linewidth}{\wd\pandoc@box}%
  \ifdim\@tempb\p@<\@tempa\p@\let\@tempa\@tempb\fi% select the smaller of both
  \ifdim\@tempa\p@<\p@\scalebox{\@tempa}{\usebox\pandoc@box}%
  \else\usebox{\pandoc@box}%
  \fi%
}
% Set default figure placement to htbp
\def\fps@figure{htbp}
\makeatother

\ifLuaTeX
  \usepackage{luacolor}
  \usepackage[soul]{lua-ul}
\else
  \usepackage{soul}
\fi


\ifLuaTeX
\usepackage[bidi=basic,shorthands=off]{babel}
\else
\usepackage[bidi=default,shorthands=off]{babel}
\fi
\ifLuaTeX
  \usepackage{selnolig} % disable illegal ligatures
\fi


\setlength{\emergencystretch}{3em} % prevent overfull lines

\providecommand{\tightlist}{%
  \setlength{\itemsep}{0pt}\setlength{\parskip}{0pt}}



 


\usepackage{booktabs}
\usepackage{caption}
\usepackage{longtable}
\usepackage{colortbl}
\usepackage{array}
\usepackage{anyfontsize}
\usepackage{multirow}
\KOMAoption{captions}{tableheading}
\makeatletter
\@ifpackageloaded{caption}{}{\usepackage{caption}}
\AtBeginDocument{%
\ifdefined\contentsname
  \renewcommand*\contentsname{Table des matières}
\else
  \newcommand\contentsname{Table des matières}
\fi
\ifdefined\listfigurename
  \renewcommand*\listfigurename{Liste des Figures}
\else
  \newcommand\listfigurename{Liste des Figures}
\fi
\ifdefined\listtablename
  \renewcommand*\listtablename{Liste des Tables}
\else
  \newcommand\listtablename{Liste des Tables}
\fi
\ifdefined\figurename
  \renewcommand*\figurename{Figure}
\else
  \newcommand\figurename{Figure}
\fi
\ifdefined\tablename
  \renewcommand*\tablename{Table}
\else
  \newcommand\tablename{Table}
\fi
}
\@ifpackageloaded{float}{}{\usepackage{float}}
\floatstyle{ruled}
\@ifundefined{c@chapter}{\newfloat{codelisting}{h}{lop}}{\newfloat{codelisting}{h}{lop}[chapter]}
\floatname{codelisting}{Listing}
\newcommand*\listoflistings{\listof{codelisting}{Liste des Listings}}
\makeatother
\makeatletter
\makeatother
\makeatletter
\@ifpackageloaded{caption}{}{\usepackage{caption}}
\@ifpackageloaded{subcaption}{}{\usepackage{subcaption}}
\makeatother
\usepackage{bookmark}
\IfFileExists{xurl.sty}{\usepackage{xurl}}{} % add URL line breaks if available
\urlstyle{same}
\hypersetup{
  pdftitle={Stratification du risque de candidémie en soins intensifs/réanimation},
  pdfauthor={BAERT Julien},
  pdflang={fr},
  colorlinks=true,
  linkcolor={blue},
  filecolor={Maroon},
  citecolor={Blue},
  urlcolor={Blue},
  pdfcreator={LaTeX via pandoc}}


\title{Stratification du risque de candidémie en soins
intensifs/réanimation}
\usepackage{etoolbox}
\makeatletter
\providecommand{\subtitle}[1]{% add subtitle to \maketitle
  \apptocmd{\@title}{\par {\large #1 \par}}{}{}
}
\makeatother
\subtitle{Mémoire Master 2 Méthodologie et statistiques en recherche
biomédicale}
\author{BAERT Julien}
\date{2026-06-04}
\begin{document}
\maketitle

\renewcommand*\contentsname{Table des matières}
{
\setcounter{tocdepth}{3}
\tableofcontents
}

\begin{center}\rule{0.5\linewidth}{0.5pt}\end{center}

\section{\texorpdfstring{\ul{Introduction}}{Introduction}}\label{introduction}

\emph{Candida} est une levure saprophyte du revêtement cutané et de la
muqueuse oro-digestive chez l'être humain. Ces champignons peuvent
provoquer des infections sanguines (candidémies) et des atteintes de
tissus profonds (candidoses invasives), représentant la pathologie
fongique nosocomiale la plus fréquente dans les pays occidentaux (1).
Les candidémies sont particulièrement fréquentes chez les patients
hospitalisés en réanimation, représentant 8\% des hémocultures positives
et 15\% des infections nosocomiales dans ces services (2). Les patients
de réanimation en présentent en effet de nombreux facteurs de risque :
certains dispositifs invasifs tels que les cathéters veineux centraux,
les cathéters de dialyse, ou la ventilation mécanique, qui sont des
outils thérapeutiques employés quasi-exclusivement en réanimation ; la
nutrition parentérale, administrée uniquement par ces cathéters profonds
; les chirurgies abdominales lourdes ou les pancréatites aigues
nécrosantes. Par ailleurs, une hospitalisation prolongée, associée à une
exposition répétée aux antibiotiques, favorisent la colonisation à
\emph{Candida}, première étape vers l'infection invasive (4,5).

Les dernières décennies ont été marquées par une incidence croissante
des candidoses invasives en réanimation, sans amélioration significative
de leur pronostic (6,7). Leur pronostic est dépendant de la précocité
d'initiation du traitement antifongique (8,9). Or le diagnostic des
candidoses invasives repose sur les hémocultures, dont la sensibilité
est imparfaite et le délai de pousse tardif, ce qui impacte négativement
le pronostic des patients (10). Pour pallier ces limites, les cliniciens
s'appuient sur les facteurs de risque et des biomarqueurs, tels que les
β-D-glucanes, pour identifier les patients à risque et introduire
précocement un traitement antifongique. Ces stratégies empiriques ou
préemptives entraînent toutefois une surutilisation des antifongiques
sans bénéfice démontré sur la survie des patients (11).

Des scores prédictifs ont été développés à partir de critères cliniques
et, dans certains cas, des variables biologiques, afin de dichotomiser
les patients en populations haut risque/bas risque (12-14). Les
recommandations actuelles portent plutôt sur une stratification des
patients en plusieurs groupes de risque, avec plusieurs niveaux de
risque (15,16). Devant les changements démographiques de patients
hospitalisés en réanimation, des changements épidémiologiques de
candidoses invasives (17) l'évolution des pratiques médicales, le
développement de nouveaux outils diagnostiques tel que la PCR, nous
proposons un nouveau score prédicitif afin de stratifier le risque de
candidémie acquise en réanimation chez les patients de réanimation en 3
catégories :

\begin{itemize}
\item
  Une population à bas risque, permettant d'éliminer l'hypothèse d'une
  candidémie
\item
  Une population à risque intermédiaire, auprès de laquelle les
  marqueurs fongiques ou les nouvelles méthodes de biologie moléculaire
  sont à proposer
\item
  Une population à haut risque (\ul{\textgreater{}} 20\%) chez il faut
  instaurer un traitement antifongique empirique d'emblée.
\end{itemize}

\section{\texorpdfstring{\ul{Méthodes}}{Méthodes}}\label{muxe9thodes}

\subsection{Source des données}\label{source-des-donnuxe9es}

Il s'agit d'une étude observationnelle monocentrique rétrospective,
utilisant les données des patients hospitalisés en soins intensifs et
réanimation au CHU de Lille depuis 2013, contenues dans l'entrepôt de
données de santé \emph{Include} du CHU de Lille. Celui-ci intègre
différentes sources informatisées de données recueillies dans le cadre
d'une prise en charge au CHU de Lille, notamment le système de gestion
des données de réanimation ICCA (Philips), le système d'information de
laboratoire MOLIS (CompuGroupMedical) et le logiciel CORA PMSI pour le
codage des diagnostics et des actes CCAM.

L'extraction des données a été réalisée par les ingénieures data
d'\emph{Include} (Mme Charlotte GEAY, Mme Léa GRUMIAUX et Mme Marie
Amélie RISPAL). La structuration des données, la comparaison des
données, les analyses statistiques ont été réalisées sur le logiciel R,
version 4.6.0.

\subsection{Population}\label{population}

La population éligible est la suivante :

\begin{itemize}
\item
  Patients hospitalisés en réanimation ou en surveillance continue au
  CHU de Lille (hôpitaux Huriez, Institut Cœur Poumon \& Salengro, hors
  service de réanimation neurochirurgicale) depuis plus de 48 heures,
\item
  Ayant été prélevés d'au moins une hémoculture après 48 heures
  d'hospitalisation en réanimation (plusieurs hémocultures peuvent être
  prélevées chez un même patient durant son séjour en réanimation).
\end{itemize}

\subsection{Critère de jugement
principal}\label{crituxe8re-de-jugement-principal}

L'objectif de notre modélisation est la stratitification du risque de
candidémie inaugurale acquise en soins intensifs/réanimation. Afin
d'être au plus proche de la réflexion du clinicien, nous procédons à un
regroupement des hémocultures par épisode de suspicion d'infection. En
effet, lors d'une suspicion d'infection, plusieurs hémocultures sont
prélevées à des temporalités proches, afin d'identifier au mieux l'agent
infectieux responsable \& un site d'infection s'il en existe (exemple :
infections de cathéters). Les hémocultures sont donc regroupées jusqu'à
ce que la dernière hémoculture prélevée pour un même patient soit
séparée d'un délai strictement supérieur à 24 heures de la prochaine
hémoculture. Une fois les hémocultures regroupées, si l'une des
hémocultures d'un groupe était positive à \emph{Candida}, nous
considérions toutes les hémocultures de ce groupe comme positives à
\emph{Candida}. Enfin, seules les données de la première hémoculture du
groupe étaient gardées, afin de placer notre modélisation à la phase la
plus initiale de l'épisode de suspicion d'infection.

L'objectif étant d'identifier une candidémie inaugurale, toutes les
hémocultures suivants une hémoculture positive à \emph{Candida} sont
supprimées.

La méthode de regroupement est résumée dans la figure
Figure~\ref{fig-regroupement}.

\begin{figure}

\centering{

\includegraphics[width=4.48in,height=\textheight,keepaspectratio]{figures/regroupement.png}

}

\caption{\label{fig-regroupement}Regroupement des hémocultures par
épisode de suspicion d'infection}

\end{figure}%

\subsection{Variables et feature
Extraction}\label{variables-et-feature-extraction}

Les variables candidates pour la modélisation étaient identifiées en
fonction de facteurs de risques identifiés dans la littérature (18). Les
données démographiques sont récoltées à l'admission et considérées comme
inchangées au décours. Ces données portaient sur le lieu
d'hospitalisation, l'âge à l'admission, les antécédents de maladie
hématologique maligne, de transplantation d'organe solide, de diabète,
de pancréatite aigue, et de tumeurs d'organes solides.

Les données cliniques et biologiques reflétants l'état à l'admission du
patient du patient sont extraites dans les 24 heures suivant l'horaire
d'admission en soins intensifs, ainsi que dans les 24 heures précédant
la réalisation de l'hémoculture : score IGS2, score SOFA, poids,
températures minimale et maximale, diurèse totale et normalisée,
exposition à des amines (adrénaline, noradrénaline, terlipressine,
dobutamine), la présence d'une ventilation invasive, la présence d'une
épuration extra-rénale, ainsi que la présence d'un remplissage
vasculaire. Les valeurs maximales de créatinémie, lactatémie et urémie
étaient récoltées ; les valeurs minimales de leucocytémie et de rapport
PaO2/FiO2. La présence d'un état de choc était définie par la présence
d'amines, d'une lactatémie \textgreater{} 2 mmol/L ou d'un score SOFA
cardiovasculaire \textgreater\{.underline\} à 3.

Afin de quantifier les expositions au cours de l'hospitalisation en
réanimation, la présence de fibroscopie(s) digestive(s), de chirurgie(s)
abdominale(s), de chirurgie(s) majeures, les durées de ventilation
mécanique invasive, de nutrition parentérale, les durées de présence de
cathéters (artériels, veineux centraux, de dialyse, ou d'ECMO), la durée
d'exposition aux antibiotiques, aux immunosuppresseurs ou aux
corticoïdes, et le nombre de culots globulaires, plaquettaires \& poches
de plasma étaient également récoltées. Les durées de neutrophilie
(définie comme \textgreater{} 7 G/L), neutropénie (définie comme
\textless{} 0,5 G/L) et de lymphopénie (définie comme lymphocytes
\textless{} 0.1 G/L) étaient identifiées.

\subsection{Analyses statistiques}\label{analyses-statistiques}

Les variables catégorielles sont décrites en effectifs \& pourcentage.
Les variables continues sont présentées par médiane, minimale et
maximale. Une description démographique \& univariée est réalisée dans
un premier temps. Sont réalisées dans un second temps des analyses
bivariées afin de de décrire le lien entre la présence d'une candidémie
\& les variables extraites. Les variables pertinentes cliniquement,
constantes avec la littérature, ou en faveur d'un sur-risque de
candidémie lors de l'analyse bivariée ont été incluses dans une analyse
multivariée. L'analyse multivariée est réalisée via régression
logistique mixte avec effet aléatoire via le patient. La modélisation
finale était évaluée par évaluation de l'aire sous la courbe ROC
(Receiving Operation Characteristic) \& par courbe de calibration.

Tous les tests statistiques étaient bilatéraux avec un niveau de
significativité fixé à p \textless{} 0.05. Les variables catégorielles
sont décrites en effectifs \& pourcentage. Les variables continues sont
présentées par médiane, minimale et maximale. \#\#\# Données manquantes

Après extraction, certaines données n'ont pu être retrouvées, faute
d'analyse ou de réalisation de l'examen. Un tableau récapitulatif des
données manquantes est présenté ci-dessous. Des imputations multiples
sont réalisées pour les données manquantes considérées \emph{Missing at
Random.}

\subsection{Cadre réglementaire}\label{cadre-ruxe9glementaire}

La conformité au Règlement Général sur la Protection des Données (RGPD)
du traitement de ces données par le CHU de Lille sur le fondement de
l'intérêt public à des fins de recherche a été validée par la Commission
Nationale de l'Informatique et des Libertés (CNIL) (Délibération
n°2019-103). Les responsables d'Include ont pris toutes les mesures pour
mener les futurs projets de recherche conformément à la loi française et
européenne.

\section{\texorpdfstring{\ul{Résultats}}{Résultats}}\label{ruxe9sultats}

\subsection{Caractéristiques de la
population}\label{caractuxe9ristiques-de-la-population}

Après extraction des hémocultures de l'entrepôt de données de santé,
nous procédons à une analyse exploratoire des données. Nous identifions
donc \textbf{104 620} hémocultures, avec \textbf{9813 séjours}, chez
\textbf{9714} individus (@fig\_flowchart). A noter qu'en moyenne, chaque
individu se voyait prélever \ul{2 hémocultures par séjour}.

\begin{figure}

\centering{

\includegraphics[width=2.2in,height=\textheight,keepaspectratio]{figures/flow_chart.png}

}

\caption{\label{fig-flowchart}Flow chart}

\end{figure}%

\subsection{Description de la
population}\label{description-de-la-population}

Dans notre population d'étude, les patients étaient majoritairement de
sexe masculin (66\%), et d'âge médian de 61 ans. 4.7\% présentaient des
antécédents de maladie hématologique maligne, 26\% de diabète, et 18\%
de tumeurs solides. A l'admission, 84\% présentaient un état de choc et
67 \% étaient ventilés de manière invasive. L'ensemble de la description
est représentée dans le @tbl\_demographique.

\begin{table}
\caption{Description démographique de la population}\tabularnewline

\fontsize{12.0pt}{14.0pt}\selectfont
\begin{tabular*}{\linewidth}{@{\extracolsep{\fill}}lc}
\toprule
\textbf{Caractéristique} & \textbf{N = 9 714}\textsuperscript{\textit{1}} \\ 
\midrule\addlinespace[2.5pt]
{\bfseries Age} & 61 (12, 97) \\ 
{\bfseries Sexe masculin} & 6 443 (66\%) \\ 
{\bfseries Antécédent de maladie hématologique maligne} & 455 (4,7\%) \\ 
{\bfseries Antécédent de diabète} & 2 514 (26\%) \\ 
{\bfseries Antécédent de pancréatite aigue} & 162 (1,7\%) \\ 
{\bfseries Antécédent de tumeur solide} & 1 740 (18\%) \\ 
{\bfseries Hospitalisation en réanimation médicale} & 7 312 (75\%) \\ 
{\bfseries IGS 2 à l'admission} & 46 (1, 135) \\ 
{\bfseries Admission pour pancréatite aigue} & 58 (0,6\%) \\ 
{\bfseries Score SOFA à l'admission} & 10,0 (0,0, 21,0) \\ 
{\bfseries Choc à l'admission} & 6 443 (84\%) \\ 
{\bfseries Poids à l'admission} & 78 (0, 482) \\ 
{\bfseries Température minimale (en °C) à l'admission} & 36,10 (18,20, 39,40) \\ 
{\bfseries Température maximale (en °C) à l'admission} & 37,90 (33,10, 42,60) \\ 
{\bfseries Diurèse normalisée} & 12 (0, 269) \\ 
{\bfseries Créatinémie maximale (en mg/L) à l'admission} & 12 (0, 251) \\ 
{\bfseries PaO2/FiO2 minimal à l'admission} & 182 (14, 648) \\ 
{\bfseries Lactatémie maximale (en mmol/L) à l'admission} & 2,10 (0,30, 30,00) \\ 
{\bfseries adm\_leuco\_min} & 11 (0, 692) \\ 
{\bfseries Neutrophiles (en G/L) à l'admission} & 11 (0, 159) \\ 
{\bfseries Lymphocytes (en G/L) à l'admission} & 0,70 (0,00, 240,70) \\ 
{\bfseries Ventilation invasive à l'admission} & 6 550 (67\%) \\ 
{\bfseries Amines à l'admission} & 5 352 (55\%) \\ 
{\bfseries Délai entre admission \& hémoculture} & 3,5 days (2,0 days, 94,9 days) \\ 
{\bfseries Décès en réanimation} & 2 746 (28\%) \\ 
\bottomrule
\end{tabular*}
\begin{minipage}{\linewidth}
\vspace{.05em}
\textsuperscript{\textit{1}}Médiane (Min, Max); n (\%)\\
\end{minipage}
\end{table}

\subsection{Modélisation}\label{moduxe9lisation}

En analyse univariée, xxx facteurs étaient associés au risque de
candidémie. Les facteurs xxx étaient au contraire associés négativement.
Les variables sélectionnées étaient ensuite analysée en régression
multivariée. Les facteurs xxx .

L'AUC étaient de xxx ; Se Sp VPP VPN. Courbe de calibration.

\subsection{Score et prédiction}\label{score-et-pruxe9diction}

AJOUTER la suite xxxxxx

\section{\texorpdfstring{\ul{Discussion}}{Discussion}}\label{discussion}

Dans cette étude monocentrique rétrospective, xxx étaient associés à un
surrisque de candidémie acquise en réanimation. Un nouveau score
prédictif était réalisé, avec des performances diagnostiques
satisfaisantes, permettant de stratifier la population décrite en 3
catégories de risque : une population à très faible risque prédit,
permettant d'éliminer presque formellement le risque de candidémie, une
population à risque modéré, pour laquelle des examens spécifiques (PCR,
marqueurs fongiques) peuvent être indiqués, et une population à haut
risque, pour laquelle il faudrait probablement introduire un traitement
antifongique empirique.

La population étudiée est fortement représentative, avec une variété de
patients hospitalisés dans différents secteurs de réanimation. En effet,
notre modélisation intéresse à la fois les secteurs médicaux,
chirurgicaux, les patients immunodéprimés ou non ; ces facteurs étant
largement employés pour stratifier le risque de candidémie. Notre étude
est également la plus grande de part sa taille (citer les autres études
qui sont toutes plus petites). Ainsi, notre étude, malgré son caractère
monocentrique, peut probablement être généralisée à une large population
; et probablement être appliquée à de multiples services de soins
intensifs \& de réanimation. La validité externe de notre score
prédictif reste à évaluer de façon prospective, à échelle
multicentrique.

Une autre force de notre étude est sa cohérence avec la littérature

Notre étude est toutefois limitée par son caractère monocentrique, qui
malgré un recrutement largerte d'informations entre les deux -
Monocentrique - Complexité de l'équation \textgreater\textgreater{}
implémenter dans logiciels au lit du malade ?

\section{\texorpdfstring{\ul{Bibliographie}}{Bibliographie}}\label{bibliographie}

1.~~~~~ Arendrup MC, Sulim S, Holm A, Nielsen L, Nielsen SD, Knudsen JD,
et al.~Diagnostic Issues, Clinical Characteristics, and Outcomes for
Patients with Fungemia.~J Clin Microbiol. sept~2011;49(9):3300‑8.~

2.~~~~~ Paiva JA, Pereira JM,~Tabah~A, Mikstacki A, De Carvalho
FB,~Koulenti~D, et al.~Characteristics and risk factors for 28-day
mortality of hospital~acquired~fungemias~in ICUs: data from the EUROBACT
study. Crit Care. 9 mars 2016;20(1):53.~

3.~~~~~ Arendrup MC. Epidemiology of invasive candidiasis: Current
Opinion in Critical Care. oct 2010;16(5):445‑52.~

4.~~~~~ Vincent JL. International Study of the Prevalence and Outcomes
of Infection in Intensive Care Units.~JAMA. 2~déc~2009;302(21):2323.~

5.~~~~~ The French Mycosis Study
Group,~Lortholary~O,~Renaudat~C,~Sitbon~K, Madec Y,~Denoeud-Ndam L, et
al.~Worrisome~trends in incidence and mortality of candidemia in
intensive care units (Paris area, 2002--2010).~Intensive Care Med. sept
2014;40(9):1303‑12~

6.~~~~~ Lamoth~F, Lockhart SR, Berkow EL, Calandra T. Changes in the
epidemiological landscape of invasive candidiasis. Journal of
Antimicrobial Chemotherapy. 1~janv~2018;73(suppl\_1):i4‑13.~

7.~~~~~ Cleveland AA, Farley MM, Harrison LH, Stein B, Hollick R,
Lockhart SR, et al.~Changes in Incidence and Antifungal Drug Resistance
in Candidemia: Results~From~Population-Based Laboratory Surveillance in
Atlanta and Baltimore, 2008-2011. Clinical Infectious Diseases.
15~nov~2012;55(10):1352‑61.~

8.~~~~~ Kollef~M, Micek S, Hampton N, Doherty JA, Kumar A. Septic Shock
Attributed to Candida Infection: Importance of Empiric Therapy and
Source Control. Clinical Infectious Diseases.
15~juin~2012;54(12):1739‑46.~

9.~~~~~ Guinea J. Global trends in the distribution of Candida species
causing candidemia. Clinical Microbiology and
Infection.~juin~2014;20:510.~

10.~ Garey KW, Rege M, Pai MP, Mingo DE, Suda KJ, Turpin RS, et al.~Time
to Initiation of Fluconazole Therapy Impacts Mortality in Patients with
Candidemia: A Multi-Institutional Study.~

11.~ Timsit JF, Azoulay E, Schwebel C, Charles PE, Cornet M, Souweine B,
et al.~Empirical Micafungin Treatment and Survival Without Invasive
Fungal Infection in Adults With ICU-Acquired
Sepsis,~\emph{Candida}~Colonization, and Multiple Organ Failure: The
EMPIRICUS Randomized Clinical Trial. JAMA. 18 oct 2016;316(15):1555.

\textbf{\_\_\_\_\_\_\_\_\_\_\_\_\_\_\_\_\_\_\_\_\_\_\_\_\_\_\_\_\_\_\_\_\_\_\_\_\_\_\_\_\_\_\_\_\_\_\_\_\_\_\_\_\_\_\_\_\_\_\_\_\_\_\_\_\_\_\_\_\_\_\_\_\_\_\_\_\_\_\_\_\_\_\_\_\_\_\_\_\_\_\_\_\_\_\_\_}

\section{\texorpdfstring{\ul{Supplementary
Data}}{Supplementary Data}}\label{supplementary-data}

\subsection{Données manquantes}\label{donnuxe9es-manquantes}

Le nombre de données manquantes à la phase à l'issue de l'extraction
sont détaillées dans @fig\_NA . Seules les variables pour lesquelles des
données manquantes étaient présentes sont représentées.

\begin{figure}[H]

{\centering \pandocbounded{\includegraphics[keepaspectratio]{M2MSR_files/figure-pdf/fig_NA-1.pdf}}

}

\caption{Données manquantes}

\end{figure}%

\textbf{\_\_\_\_\_\_\_\_\_\_\_\_\_\_\_\_\_\_\_\_\_\_\_\_\_\_\_\_\_\_\_\_\_\_\_\_\_\_\_\_\_\_\_\_\_\_\_\_\_\_\_\_\_\_\_\_\_\_\_\_\_\_\_\_\_\_\_\_\_\_\_\_\_\_\_\_\_\_\_\_\_\_\_\_\_\_\_\_\_\_\_\_\_\_\_\_}

\subsection{Comparaisons bivariées}\label{comparaisons-bivariuxe9es}

Nous réalisons une comparaison bi-variées, avec correction du risque
d'inflation du risque alpha via la méthode de \emph{false detection rate
(FDR)}.

\begin{table}
\caption{Analyse bivariée de la survenue d'une candidémie}\tabularnewline

\fontsize{12.0pt}{14.0pt}\selectfont
\begin{tabular*}{\linewidth}{@{\extracolsep{\fill}}lcccc}
\toprule
\textbf{Caractéristique} & \textbf{Négative}  N = 29 577\textsuperscript{\textit{1}} & \textbf{Positive}  N = 376\textsuperscript{\textit{1}} & \textbf{p-valeur}\textsuperscript{\textit{2}} & \textbf{q-valeur}\textsuperscript{\textit{3}} \\ 
\midrule\addlinespace[2.5pt]
\multicolumn{5}{l}{{\bfseries Données démographiques}} \\[2.5pt] 
\midrule\addlinespace[2.5pt]
Hospitalisation en réanimation médicale & 21 341 (72\%) & 290 (77\%) & 0,032 & 0,15 \\ 
Age & 60 (12, 97) & 59 (16, 88) & 0,7 & 0,7 \\ 
Sexe & 20 189 (68\%) & 253 (67\%) & 0,7 & 0,7 \\ 
Antécédent de maladie hématologique maligne & 1 310 (4,5\%) & 25 (6,7\%) & 0,037 & 0,15 \\ 
Antécédent de diabète & 7 598 (26\%) & 103 (28\%) & 0,4 & 0,6 \\ 
Antécédent de pancréatite aigue & 492 (1,7\%) & 10 (2,7\%) & 0,13 & 0,4 \\ 
Antécédent de tumeur solide & 4 670 (16\%) & 67 (18\%) & 0,3 & 0,6 \\ 
Antécédent de transplantation d'organe solide & 1 647 (5,6\%) & 17 (4,5\%) & 0,4 & 0,6 \\ 
\midrule\addlinespace[2.5pt]
\multicolumn{5}{l}{{\bfseries Données à l'admission}} \\[2.5pt] 
\midrule\addlinespace[2.5pt]
Etat de choc & 20 343 (86\%) & 283 (90\%) & 0,018 & 0,030 \\ 
IGS 2 & 47 (0, 135) & 52 (8, 135) & <0,001 & <0,001 \\ 
Pancréatite aigue & 295 (1,1\%) & 2 (0,6\%) & 0,6 & 0,6 \\ 
Poids & 80 (0, 482) & 80 (0, 250) & 0,5 & 0,5 \\ 
Température minimale (en °C) & 36,10 (18,20, 39,40) & 36,10 (28,80, 39,00) & 0,2 & 0,3 \\ 
Température maximale (en °C) & 37,90 (33,10, 42,60) & 37,85 (34,10, 40,70) & 0,2 & 0,3 \\ 
Diurèse totale & 900 (0, 12 600) & 700 (0, 5 005) & <0,001 & <0,001 \\ 
Diurèse normalisée & 11 (0, 269) & 9 (0, 62) & <0,001 & <0,001 \\ 
Créatinémie maximale (en mg/L) & 12 (0, 251) & 16 (1, 178) & <0,001 & <0,001 \\ 
Urée maximale (en g/L) & 0,56 (0,03, 7,91) & 0,70 (0,03, 4,20) & <0,001 & <0,001 \\ 
PaO2/FiO2 minimal & 164 (14, 648) & 142 (26, 543) & 0,003 & 0,006 \\ 
Lactatémie maximale (en mmol/L) & 2,30 (0,30, 30,00) & 2,80 (0,50, 23,00) & <0,001 & <0,001 \\ 
Lactatémie moyenne (en mmol/L)  & 1,75 (0,24, 22,00) & 2,15 (0,42, 16,60) & <0,001 & <0,001 \\ 
Leucocytes (en G/L) & 11 (0, 692) & 11 (0, 389) & 0,5 & 0,6 \\ 
Neutrophiles (en G/L) & 10 (0, 159) & 10 (0, 70) & 0,8 & 0,8 \\ 
Lymphocytes (en G/L) & 0,70 (0,00, 240,70) & 0,60 (0,00, 191,20) & 0,034 & 0,049 \\ 
Ventilation mécanique invasive & 21 555 (73\%) & 283 (75\%) & 0,3 & 0,4 \\ 
Dialyse & 2 903 (9,8\%) & 50 (13\%) & 0,024 & 0,037 \\ 
Culot Globulaires & 2 997 (10\%) & 68 (18\%) & <0,001 & <0,001 \\ 
Plasma frais congelé & 732 (2,5\%) & 22 (5,9\%) & <0,001 & <0,001 \\ 
Culot plaquettaire & 1 371 (4,6\%) & 42 (11\%) & <0,001 & <0,001 \\ 
Transfusion & 3 816 (13\%) & 94 (25\%) & <0,001 & <0,001 \\ 
Amines  & 17 543 (59\%) & 249 (66\%) & 0,007 & 0,012 \\ 
\midrule\addlinespace[2.5pt]
\multicolumn{5}{l}{{\bfseries Données dans les 48 heures précédant la suspicion d'infection}} \\[2.5pt] 
\midrule\addlinespace[2.5pt]
Délai entre admission et prélèvements de l'hémoculture & 9 days (2 days, 520 days) & 10 days (2 days, 293 days) & 0,5 & 0,5 \\ 
Etat de choc & 11 775 (60\%) & 218 (76\%) & <0,001 & <0,001 \\ 
Température minimale  & 36,70 (15,10, 40,00) & 36,60 (32,50, 39,00) & 0,016 & 0,020 \\ 
Température maximale  & 38,30 (33,80, 42,20) & 38,20 (35,40, 40,30) & 0,4 & 0,4 \\ 
Diurèse totale  & 1 100 (0, 19 300) & 650 (0, 7 200) & <0,001 & <0,001 \\ 
Diurèse normalisée  & 14 (0, 243) & 9 (0, 67) & <0,001 & <0,001 \\ 
PaO2/FiO2  & 200 (16, 648) & 178 (31, 524) & 0,048 & 0,056 \\ 
Créatinémie  & 9 (1, 148) & 12 (1, 92) & <0,001 & <0,001 \\ 
Urémie  & 0,57 (0,02, 4,50) & 0,71 (0,06, 4,49) & <0,001 & <0,001 \\ 
Lactatémie maximale & 1,40 (0,10, 27,00) & 1,80 (0,50, 20,00) & <0,001 & <0,001 \\ 
Lactatémie moyenne  & 1,23 (0,10, 20,00) & 1,60 (0,47, 15,83) & <0,001 & <0,001 \\ 
Neutropénie min  & 10 (0, 81) & 12 (0, 50) & 0,010 & 0,012 \\ 
Leucocytes (en G/L) & 12 (0, 580) & 12 (0, 125) & 0,7 & 0,7 \\ 
Lymphopénie  & 0,90 (0,00, 465,30) & 0,70 (0,00, 3,80) & <0,001 & <0,001 \\ 
Glucanes positifs & 6 623 (39\%) & 143 (57\%) & <0,001 & <0,001 \\ 
Mannanes positifs & 3 749 (23\%) & 104 (44\%) & <0,001 & <0,001 \\ 
Ventilation invasive  & 19 850 (67\%) & 295 (78\%) & <0,001 & <0,001 \\ 
Dialyse  & 4 274 (14\%) & 100 (27\%) & <0,001 & <0,001 \\ 
KTA  & 18 227 (62\%) & 266 (71\%) & <0,001 & <0,001 \\ 
VVC  & 22 787 (77\%) & 326 (87\%) & <0,001 & <0,001 \\ 
Cathéter de dialyse  & 5 864 (20\%) & 119 (32\%) & <0,001 & <0,001 \\ 
ECMO  & 1 563 (5,3\%) & 24 (6,4\%) & 0,3 & 0,4 \\ 
CGR  & 2 237 (7,6\%) & 64 (17\%) & <0,001 & <0,001 \\ 
PFC  & 147 (0,5\%) & 7 (1,9\%) & 0,003 & 0,005 \\ 
CP  & 1 049 (3,5\%) & 47 (13\%) & <0,001 & <0,001 \\ 
Transfusion  & 2 855 (9,7\%) & 84 (22\%) & <0,001 & <0,001 \\ 
Amines  & 9 802 (33\%) & 177 (47\%) & <0,001 & <0,001 \\ 
\midrule\addlinespace[2.5pt]
\multicolumn{5}{l}{{\bfseries Données au cours de l'hospitalisation préalablement à l'hémoculture}} \\[2.5pt] 
\midrule\addlinespace[2.5pt]
Durée de ventilation invasive  & 130 (0, 6 205) & 158 (0, 6 245) & 0,030 & 0,045 \\ 
Durée de nutrition parentérale  & 0,0 (0,0, 161,0) & 1,5 (0,0, 77,0) & <0,001 & <0,001 \\ 
Durée de VVC  & 9 (0, 235) & 10 (0, 218) & 0,030 & 0,045 \\ 
Durée de KTA  & 8 (0, 178) & 9 (0, 93) & 0,026 & 0,045 \\ 
Durée de KTD  & 0 (0, 227) & 0 (0, 144) & <0,001 & <0,001 \\ 
Durée ECMO  & 0,00 (0,00, 88,00) & 0,00 (0,00, 90,00) & 0,14 & 0,2 \\ 
Durée ATBT  & 8 (0, 244) & 8 (0, 151) & 0,021 & 0,043 \\ 
Durée CTC  & 0 (0, 267) & 2 (0, 270) & <0,001 & <0,001 \\ 
Durée immunosuppression  & 0,00 (0,00, 163,00) & 0,00 (0,00, 35,00) & 0,5 & 0,5 \\ 
Durée neutropénie  & 0,00 (0,00, 23,00) & 0,00 (0,00, 16,00) & 0,2 & 0,3 \\ 
Durée neutrophilie  & 3,0 (0,0, 69,0) & 4,0 (0,0, 69,0) & 0,003 & 0,007 \\ 
Durée lymphopénie  & 0,00 (0,00, 59,00) & 0,00 (0,00, 14,00) & <0,001 & <0,001 \\ 
Nombres CGR avant l'hémoculture & 1 (0, 132) & 3 (0, 85) & <0,001 & <0,001 \\ 
Nombres PFC avant l'hémoculture & 0,00 (0,00, 326,00) & 0,00 (0,00, 80,00) & <0,001 & <0,001 \\ 
Nombres CP avant l'hémoculture & 0,00 (0,00, 104,00) & 0,00 (0,00, 48,00) & <0,001 & <0,001 \\ 
Fibroscopie pendant l'hospitalisation & 5 291 (18\%) & 78 (21\%) & 0,2 & 0,2 \\ 
Chirurgie majeure pendant l'hospitalisation & 8 599 (29\%) & 120 (32\%) & 0,2 & 0,3 \\ 
Chirurgie abdominale pendant l'hospitalisation & 5 394 (18\%) & 68 (18\%) & >0,9 & >0,9 \\ 
\bottomrule
\end{tabular*}
\begin{minipage}{\linewidth}
\vspace{.05em}
\textsuperscript{\textit{1}}n (\%); Médiane (Min, Max)\\
\textsuperscript{\textit{2}}test du khi-deux d'indépendance; test de Wilcoxon-Mann-Whitney\\
\textsuperscript{\textit{3}}Correction du risque global d'erreur pour tests multiples\\
\end{minipage}
\end{table}

\subsection{Modélisation stepwise à partir d'une imputation
simple}\label{moduxe9lisation-stepwise-uxe0-partir-dune-imputation-simple}

A titre exploratoire, des modélisations stepwise forward et backward ont
été réalisées à partir d'un unique jeu imputé.

\begin{figure}

\begin{minipage}[t]{0.50\linewidth}

\begin{figure}[H]

{\centering \includegraphics[width=6.66in,height=\textheight,keepaspectratio]{figures/FP_fwd.png}

}

\subcaption{Forrest Plot de la régression forward}

\end{figure}%

\end{minipage}%
%
\begin{minipage}[t]{0.50\linewidth}

\begin{figure}[H]

{\centering \includegraphics[width=6.66in,height=\textheight,keepaspectratio]{figures/FP_bwd.png}

}

\subcaption{Forrest Plot de la régression backward}

\end{figure}%

\end{minipage}%
\newline
\begin{minipage}[t]{0.50\linewidth}

\begin{figure}[H]

{\centering \includegraphics[width=6.66in,height=\textheight,keepaspectratio]{figures/ROC_fwd.png}

}

\subcaption{ROC curve de la régression forward}

\end{figure}%

\end{minipage}%
%
\begin{minipage}[t]{0.50\linewidth}

\begin{figure}[H]

{\centering \includegraphics[width=6.66in,height=\textheight,keepaspectratio]{figures/ROC_bwd.png}

}

\subcaption{ROC curve de la régression backward}

\end{figure}%

\end{minipage}%
\newline
\begin{minipage}[t]{0.50\linewidth}

\begin{figure}[H]

{\centering \includegraphics[width=3.2in,height=\textheight,keepaspectratio]{figures/CC_fwd.png}

}

\subcaption{Courbe de calibration de la régression forward}

\end{figure}%

\end{minipage}%
%
\begin{minipage}[t]{0.50\linewidth}

\begin{figure}[H]

{\centering \includegraphics[width=3.2in,height=\textheight,keepaspectratio]{figures/CC_bwd.png}

}

\subcaption{Courbe de calibration de la régression backward}

\end{figure}%

\end{minipage}%
\newline
\begin{minipage}[t]{0.50\linewidth}
Forrest plots, AUC \& Courbes de calibration des modélisations stepwise
forward \& backward après imputation simple\end{minipage}%

\end{figure}%

\subsection{Code}\label{code}

\href{https://github.com/djoubar/M2MSR_candida}{Voir le code sur GitHub}




\end{document}
